\section{Introduction}
\label{sec:introduction}

A mass gap and an isolated particle mass are different spectral statements.
A positive spectral measure may start at a strictly positive threshold and
still have no atoms. For the spin field of the two-dimensional $O(3)$ model,
we prove that the lowest mass is an atom separated from all remaining mass
support. The proof uses the massive continuum construction in
\cite{OpenAI-O3}, including its uniform finite-volume estimates.

\subsection{The field and the result}

Let $\phi=(\phi^1,\phi^2,\phi^3)$ be the Hermitian vector field obtained
from the nearest-neighbor $O(3)$ model by the fixed lattice-spacing and
field-normalization prescription of \cite[Theorem~1.2]{OpenAI-O3}.
Here the Euclidean lattice is two-dimensional and the reconstructed
spacetime has one time and one space dimension. Internal $O(3)$ symmetry
and the K\"all\'en--Lehmann representation give
\begin{equation}
 W_2^{ij}(x)=\delta_{ij}\int_{[0,\infty)}
       \Delta_+(x;\mu^2)\,\rho(d\mu^2),
 \label{intro:spectral}
\end{equation}
where $\Delta_+(x;\mu^2)$ is the positive-frequency free scalar
two-point distribution and $\rho$ is a positive measure. The variable
$\mu$ denotes mass; the argument of $\rho$ is mass squared. This
representation records the masses to which the field couples
\cite{Kallen1952,Lehmann1954}.

\begin{theorem}[An isolated lowest mass]
\label{intro:main}
For the field $\phi$ just specified, there exist $m_1,Z,\delta>0$ and a
positive measure $\rho_{\mathrm{rest}}$ such that
\begin{equation}
 \rho=Z\delta_{m_1^2}+\rho_{\mathrm{rest}},
 \qquad
 \supp\rho_{\mathrm{rest}}\subset[(m_1+\delta)^2,\infty).
 \label{intro:atom}
\end{equation}
In particular, $m_1$ is the smallest mass in the support of the spin-field
two-point function.
\end{theorem}

The atom in \eqref{intro:atom} is the isolated particle contribution
to the two-point function. Both its positive weight and its separation
from the remaining spectrum must be established: neither follows from
exponential decay alone. Our use of the continuum construction is made
precise in Section~\ref{sec:inputs}. The new argument begins with its
uniform torus estimates and does not assume that an eigenvalue at finite
circumference persists in the infinite-volume limit.

\subsection{Context and methods}

The two-dimensional nonlinear sigma models are basic examples in which
short-distance behavior and massive long-distance behavior must be
reconciled. Polyakov's work on the interaction of Goldstone particles
identified the role of asymptotic freedom in these models
\cite{Polyakov1975}. The factorized scattering description of the
$O(N)$ sigma model developed by Zamolodchikov and Zamolodchikov gives
a particle-based account of its expected massive continuum physics
\cite{Zamolodchikov1979}. Relating such a description to the field
obtained from a specified lattice limit requires control of that limit
and of the field's spectral measure. Within the massive-particle
description, Balog and Niedermaier computed spin-field form factors and
spectral densities \cite{BalogNiedermaier1997}. The present argument
establishes the isolated mass directly for the field obtained from the
lattice prescription.

In constructive field theory, Glimm, Jaffe, and Spencer proved an
isolated particle mass for weakly coupled $P(\phi)_2$ models by a
particle cluster expansion \cite{GlimmJaffeSpencer1974}. The setting
and the estimates used here are different: our starting point is the
finite-volume control accompanying the $O(3)$ lattice limit.

The companion construction \cite{OpenAI-O3} provides the continuum
field, its Osterwalder--Schrader reconstruction
\cite{OsterwalderSchrader1973,OsterwalderSchrader1975}, a gap above the
vacuum, and a nonzero one-field state. It also supplies two inputs that
are decisive here: exponentially small partition-function defects under
doubling of large physical squares, uniformly in the lattice cutoff,
and uniform factorial moment bounds for normalized local spin sums.
We use these quantitative estimates to extract the isolated mass.

Several classical correlation tools enter the passage from transfer
energies to field correlations. The positivity arguments belong to the
Griffiths--Ginibre theory and its extensions to multicomponent rotators
\cite{Ginibre1970,Dunlop1976}. We give the particular three-component
argument in full, because it must control arbitrary spin polynomials
invariant under simultaneous spin inversion. Positive association
\cite{FKG1971} and the
Edwards--Sokal coupling \cite{EdwardsSokal1988} relate spin correlations
to connection events for embedded Ising signs.
The crossing theorem of K\"ohler-Schindler and Tassion
\cite{KohlerSchindlerTassion2023} supplies circuits using only symmetry
and positive association. These circuits attach a definite amount of
normalized spin-field weight to connections; the normalization is
essential as the lattice spacing tends to zero.

The use of bounded-energy state counts to constrain particle content is
related to the phase-space criteria of Haag and Swieca
\cite{HaagSwieca1965} and Buchholz and Wichmann
\cite{BuchholzWichmann1986}. Here the needed estimate is obtained directly
from finite-volume transfer traces. Testing an interval of momenta then
constrains the mass support without presupposing that its points are atoms.

\subsection{Proof strategy}

Write $m_*$ for the bottom of the mass support of $\rho$; the source
gap and nonzero one-field norm imply $0<m_*<\infty$. At lattice cutoff
$N$ and spatial circumference $w$, let $e_N(w)$ be the least positive
transfer energy, and let $\mathcal N_{N,w}(u)$ count all transfer
energies at most $u$, with multiplicity and with the vacuum included.
The proof first compares $e_N(w)$ with $m_*$, then bounds the number
of low transfer energies and uses that bound to constrain the mass support.

\begin{enumerate}
\item Section~\ref{sec:transfer} derives rectangular trace bounds and
local torus-to-plane comparison from the square partition-function
input. This also gives a positive lower bound for $e_N(w)$, uniform
in the cutoff and in sufficiently large circumferences.

\item Section~\ref{sec:parity} bounds covariances of even spin
polynomials by the square of a two-point correlation. Consequently,
among row functions unchanged by simultaneous spin inversion, every
nonvacuum transfer energy is at least $2e_N(w)$.
The lowest excitation is odd, with a quantitative overlap
that allows its sign to be tested by spin clusters.

\item Section~\ref{sec:visibility} combines those clusters with local
circuits to prove
\[
 \liminf_{w\to\infty}\liminf_{N\to\infty}e_N(w)\ge m_*.
\]
Here and in the counting step, $w$ runs through a fixed dyadic sequence
of physical circumferences. This comparison identifies the mass scale
needed for counting, without selecting a limiting eigenvector.

\item Section~\ref{sec:counting} proves
\[
 \limsup_{N\to\infty}\mathcal N_{N,w}(3m_*/2)\le Cw.
\]
Exchanging the directions of a torus compares the heat trace at
circumference $2w$ with the square of the trace at circumference $w$.
A positive polynomial approximation to an energy cutoff shows that,
below the two-excitation threshold, doubling the circumference can at
most double the count, up to summable errors.

\item Section~\ref{sec:mass-support} uses spatial momentum projections
and convergence of smeared two-point functions to show that $J$
distinct masses near $m_*$ require at least $cJw$ such states. The
constant $c>0$ is independent of $J$. Hence there are only finitely
many masses near $m_*$, and their smallest point is the atom in
Theorem~\ref{intro:main}.
\end{enumerate}

Throughout, cutoff limits are taken at fixed physical circumference
before that circumference tends to infinity. The trace and moment
estimates permit the comparisons needed in this order.
